\section{Experimental setup}
\label{sec:setup}

\subsection{Detectors}
We evaluated \NDetectors detectors (Table~\ref{tab:inventory}): the most downloaded Hugging Face text classifiers labeled as prompt-injection detectors that output an injection decision, plus the agent-oriented detectors released in 2026. Six are gated behind a license agreement, which we accepted: Prompt Guard~1 and Prompt Guard~2 (86M and 22M)~\cite{llamafirewall}, ProtectAI's small v2 model, and two Sentinel models. The set spans three generations: DeBERTa-v3 classifiers from 2023--2025 (ProtectAI, JasperLS, TestSavant, PIGuard~\cite{piguard}, and Prompt Guard), ModernBERT- and Qwen3-based classifiers from 2025--2026 (the two Sentinels, and Sheltron, Wolf Defender, and Horizon-Labs, whose cards target agent inputs), and a fine-tuned 1.5B generative classifier (Prismor). Each detector was run as its model card documents, including its context length (512 to 2{,}048 tokens), its segmentation of long inputs, its input template, and, for Sheltron, its role and source tokens; for Prompt Guard~1 we used Meta's indirect-injection score, the sum of its injection and jailbreak probabilities. \ref{app:configs} gives their Hugging Face identifiers, revisions, and input settings.

InjecGuard returned the same scores as PIGuard on every input, so we report it once. Sheltron's encoder also accepts a trusted ``user goal'' segment; this use is undocumented, but we report it as a task-aware variant because no other detector can take the user's request as input. Two detectors, \texttt{deepset/deberta-v3-base-injection} and \texttt{fmops/distilbert-prompt-injection}, flagged every benign README chunk and were omitted (\ref{app:excluded}); one gated detector (CodeIntegrity) did not grant access.

\begin{table}[t]
\centering\small
\caption{Detectors evaluated, by release year.}
\label{tab:inventory}
% generated by analysis/make_tables.py
\begin{tabular*}{\columnwidth}{@{\extracolsep{\fill}}llrl@{}}
\toprule
Detector & Year & Parameters & Agent data$^{\mathrm{a}}$ \\
\midrule
JasperLS & 2023 & 184M & No \\
ProtectAI v1 & 2023 & 184M & No \\
Prompt Guard 1 & 2024 & 279M & No \\
ProtectAI v2-small & 2024 & 142M & No \\
ProtectAI v2 & 2024 & 184M & No \\
TestSavant-L & 2024 & 184M & No \\
PIGuard & 2025 & 184M & No \\
Prompt Guard 2 (22M) & 2025 & 71M & No \\
Prompt Guard 2 (86M) & 2025 & 279M & No \\
Qualifire Sentinel & 2025 & 396M & No \\
Sentinel v2 & 2025 & 596M & No \\
Horizon-Labs & 2026 & 308M & Yes \\
Prismor-1.5B & 2026 & 1.5B & No \\
Sheltron & 2026 & 68M & Yes \\
Wolf Defender & 2026 & 308M & Eval.$^{\mathrm{b}}$ \\
\bottomrule
\end{tabular*}

\par\smallskip\footnotesize\raggedright
$^{\mathrm{a}}$The model card states training on agent or tool-output inputs.
$^{\mathrm{b}}$The card reports evaluation on AgentDojo-derived data; training data are not public.
\end{table}

\subsection{Task-aware judges}
To test whether information that the detectors lack would change the outcome, we gave an instruction-tuned model the user's request and one tool output and asked whether the output tries to make the assistant do something the user did not ask for (\ref{app:prompt}). The score is the probability of ``Yes'' relative to ``No'' at the first generated token, summed over token variants. We used Qwen2.5-7B-Instruct~\cite{qwen25} and Llama-3.1-8B-Instruct~\cite{llama3}, both quantized to 4 bits (NF4). The Llama model's training data end in December 2023, six months before AgentDojo was released, so it cannot have seen AgentDojo's tasks or attack templates. The 4-bit Qwen judge matches the same model served in bf16 to within 0.02 AUROC on AgentDojo.

\subsection{Datasets and splits}
Table~\ref{tab:data} summarizes the evaluation sets of Section~\ref{sec:method}: general benign text, two agent benchmarks, and one injection benchmark.

\begin{table}[t]
\centering\small
\caption{Evaluation sets (AD: AgentDojo v1.2; TB: $\tau$-bench). AD-Matched and AD-Shift use different AgentDojo attack templates; the benign side of TB-Matched is TB-Clean.}
\label{tab:data}
\begin{tabular*}{\columnwidth}{@{\extracolsep{\fill}}lrrl@{}}
\toprule
Set & Benign (n) & Injected (n) & Content \\
\midrule
General text & \NGeneral & -- & Prose, docs \\
AD-Clean & \NCleanOutputs & -- & No injection \\
AD-Matched & \NADBenign & \NADInj & 3 templates \\
AD-Shift & \NADSBenign & \NADSInj & 5 templates \\
TB-Matched & \NTBBenign & \NTBInj & JSON outputs \\
BIPIA & \NBIBenign & \NBIInj & Email, table, code \\
\bottomrule
\end{tabular*}
\end{table}

\subsection{Hardware configuration}
Detectors ran in PyTorch~2.8 with Hugging Face Transformers~4.57 (Transformers~5 for Wolf Defender), and judges in 4-bit NF4 through bitsandbytes. The main scores were computed on a 20\,GB partition of an NVIDIA H100 GPU; the experiments of Sections~\ref{sec:robust} and~\ref{sec:cost} ran on NVIDIA A40 GPUs of a university cluster. Scoring all sets takes minutes per detector and a few hours per judge.

\subsection{Reproducibility}
\label{sec:impl}
Every evaluation set is rebuilt by scripts from the public benchmarks (AgentDojo~0.1.34 with benchmark version v1.2, and the BIPIA, $\tau$-bench, and InjecAgent repositories), so no benchmark data are redistributed, and every number, table, and figure is generated by scripts from stored detector scores. For Sections~\ref{sec:robust} and~\ref{sec:community}, we rebuilt the AgentDojo sets on a second machine. They have the same size and label sequence as the originals, and rescoring their benign outputs with PIGuard and Prompt Guard~2 (86M) on a different GPU reproduced the stored scores to within \ReproMaxDiff. Injected outputs can differ in wording, because the original build assigned injection tasks with Python's per-process string hash (the released scripts now fix the seed); on the rebuilt set, Prompt Guard~2 (86M) catches \PGtRebuiltTpr of injections at 1\% FPR (stored: \MPGtAD) and PIGuard \PGRebuiltTpr (stored: \ADPGTpr).
